% TOP_PROBLEM: {"id": 3249, "title": "Choosability of Graph Powers", "category_id": 3, "category": {"id": 3, "name": "graph_theory", "display_name": "Graph Theory", "description": "Problems involving graphs, networks, and their properties.", "slug": "graph-theory", "order_index": 3, "created_at": "2026-07-31T15:26:25.671Z"}, "status": "open", "problem_number": "OPG-53019", "set_id": 12}
% FIRST_POSTED: 2026-10-01
% ENTRY_KIND: statement_only
% UnsolvedMath ID 3249; source catalogue code OPG-53019
% !TeX program = lualatex
% Definition and sources only; this file is not an LLM solution attempt.
\documentclass[11pt,a4paper]{article}
\usepackage[margin=25mm]{geometry}
\usepackage{fontspec}
\setmainfont{lmroman10-regular.otf}[BoldFont=lmroman10-bold.otf,
 ItalicFont=lmroman10-italic.otf,BoldItalicFont=lmroman10-bolditalic.otf]
\usepackage{amsmath,amssymb,mathtools}
\usepackage{unicode-math}
\setmathfont{latinmodern-math.otf}
\AtBeginDocument{\let\setminus\smallsetminus}
\usepackage{xurl,hyperref}
\hypersetup{colorlinks=true,urlcolor=blue}
\setcounter{secnumdepth}{0}
\setlength{\parindent}{0pt}
\setlength{\parskip}{5pt}
\setlength{\emergencystretch}{3em}
\begin{document}
\section*{Choosability of Graph Powers}\label{3249}
{\small UnsolvedMath ID 3249 · OPG-53019 · Graph Theory · OpenGarden}
\subsection{Definitions and mathematical statement}
Let $G$ be a finite nonempty simple graph. Its square $G^2$ joins two
distinct vertices when a path of length one or two connects them in $G$.
For any graph $H$, let $\chi(H)$ be its ordinary chromatic number and
$\operatorname{ch}(H)$ its choice number. The latter is the least $q$
such that every assignment of $q$ colours to each vertex permits a proper
colouring using one colour from each vertex's list.

The problem asks whether there exists a function
$f:\mathbb N_{\ge1}\to\mathbb R_{\ge0}$ satisfying
\[
 \lim_{k\to\infty}\frac{f(k)}{k^2}=0,
 \qquad
 \operatorname{ch}(G^2)\le f(\chi(G^2))
 \quad\text{for every such }G.
\]
The same function must work for all graphs, regardless of their number
of vertices. No monotonicity or effective computability of $f$ is assumed.
Equivalently, for every $\varepsilon>0$ there should exist $K$ such that
$\operatorname{ch}(G^2)\le\varepsilon\chi(G^2)^2$ whenever
$\chi(G^2)\ge K$.

The normalization is natural: if $d=\Delta(G)$, then a vertex and its
neighbours form a clique in $G^2$, while every square-degree is at most
$d^2$. Consequently the elementary greedy bound gives
$\operatorname{ch}(G^2)\le d^2+1\le\chi(G^2)^2$.
The requested improvement is uniform and asymptotically smaller than
quadratic; it need not make list chromatic number equal ordinary
chromatic number. The exponent two denotes graph distance as well as
the separate quadratic growth benchmark, in their respective positions.
\subsection{Short English statement}
Is there a universal subquadratic upper bound for the choice number of a graph square in terms of that square’s chromatic number?
\subsection{Sources}
\begin{itemize}
\item[S1] ULAM.AI and the UnsolvedMath Contributors, supplied problems.json, record 3249 (OPG-53019), OpenGarden collection. Catalogue identity and source transcription; CC BY 4.0.
\url{https://huggingface.co/datasets/ulamai/UnsolvedMath}.
\item[S2] Open Problem Garden, Choosability of Graph Powers. Original problem and discussion; the mathematical formulation below is expanded from the supplied source record.
\url{http://www.openproblemgarden.org/op/choosability_of_graph_powers}.
\end{itemize}
\end{document}
